\begin{titlepage}
	\makeatletter
	\begin{center}
		\vspace*{5cm}
		\@title
		\vfill
		{\small Last Updated on \@date}
	\end{center}
	\makeatother
\end{titlepage}

%\maketitle

\chapter{Preface}
Disclaimer, this was written by physicists. We'd love to include advice for other subjects. If anyone is willing to contribute or help us out with adding content for any of the other subjects please reach out to us at \href{mailto:vishal.vetrivel@cbs.ac.in}{vishal.vetrivel@cbs.ac.in} or \href{mailto:dhiyanesh.gowrisankar@cbs.ac.in}{dhiyanesh.gowrisankar@cbs.ac.in} or if this no longer works you can contact the science club at \href{mailto:science.club@cbs.ac.in}{science.club@cbs.ac.in} and they'll pass it on to us.

A much more efficient way to contribute to this document would be to open pull requests on the \href{https://github.com/VishalVSV/ReportGuide}{github repository}\footnote{\url{https://github.com/VishalVSV/ReportGuide}, if for some reason that link is difficult to click.}. You can also build this document yourself from the source.

This document is, by nature, incomplete and will require extensive contributions from students in many fields and subjects, to reach a state where it can be called complete. As such we invite everyone to send us any valuable information or help us improve this document directly (pointing out typos and grammatical mistakes also count as help). We hope that this will be a community driven project that matures over time to help many students.

We also wrote most of this document from the perspective of someone using a fully featured \LaTeX\ editor. More specifically we write this from the perspective of TeXstudio (\url{https://texstudio.org/}). While we do recommend it for daily use, it is not necessary nor is it the most popular tool for the job. TeXstudio is also open source and uses the GPL v2 license which makes it free software.

%\section*{Acknowledgements}

\section*{Useful Resources}
\begin{itemize}
	\item \url{https://www.overleaf.com/learn}
	\item \url{https://tex.stackexchange.com/}
\end{itemize}
\tableofcontents
